\par\noindent\textbf{(1)} 二能带本征频率E±(k)=±|v+we\textasciicircum{}\{ika\}|=±√(v²+w²+2vwcoska)。最小直接能隙为2|v−w|，在v=w且ka=π处关闭；只要v≠w，能隙打开且手征算符Γ=diag(1,−1)满足\{Γ,H\}=0。
\par\noindent\textbf{(2)} 令q(k)=v+we\textasciicircum{}\{ika\}，随k穿过第一布里渊区在复平面描绘以v为圆心、w为半径的圆。绕原点的整数绕数ν=(1/2π)∫∂karg q dk，在w>v时为1，在v>w时为0；v=w圆穿过原点，拓扑数因能隙关闭而无法定义。
\par\noindent\textbf{(3)} 取从A1端开始的半无限开链，零能且只在A子格非零的态满足v A\_n+w A\_\{n+1\}=0，故A\_n=A1(−v/w)\textasciicircum{}\{n−1\}。只有|v/w|<1时波函数可归一，归一系数√[1−(v/w)²]，端态局域长度ξ≈a/ln(w/v)。
\par\noindent\textbf{(4)} 有限长链左右两端尾部重叠使原本零能态劈裂量级∝(v/w)\textasciicircum{}N，N为晶胞数；若加入不同A/B位能会破坏手征对称并把零能推离零，但保留能隙的小扰动仍可能有局域态。拓扑边界态数还依端部切断位置，不是单由体内w>v就能不加说明确定。
